\section{Introduction}
\label{sec:intro}

Center manifold theory provides an elegant and powerful tool for the analysis of dynamics of systems in the vicinity of the bifurcation point \citep{carr82,carr83b,carr83,sijbrand85,guckenheimer83,wiggins03,haragus11}. In such systems the essential dynamics are primarily driven by the dynamics in an invariant manifold that is a continuation of the center subspace of the linearized system into the nonlinear regime. The center manifold has the same dimension as the center space leading to a low-dimensional representation of what might originally be a very high dimensional system. The ideas underlying the theory are powerful and extend beyond center manifolds, having found applicability to general invariant manifolds \citep{roberts89,roberts14,haller16}. In the hydrodynamic literature a similar approach has proliferated under the general framework of ``weakly nonlinear analysis'' \citep{sipp07,cross09,meliga11,meliga12}, although analytical work utilizing ideas of scale separation have been undertaken much earlier as well \citep{stuart60,watson60,newell69}.  \citet{fujimara91} has claimed that the two methods (center manifold and weakly nonlinear analysis) are effectively equivalent although, the underlying ideas of the two methods are somewhat distinct. 

While the importance of invariant manifold in general, and center manifold in particular, has long been recognized, its application was initially largely confined to small ordinary differential equation type systems \citep{carr82,wiggins03,guckenheimer83}. There has however been a steady progress towards application of invariant manifold theory to larger and more complex systems. The work of \citet{coullet83} is seminal in this regard which presented a method to directly obtain the normal form representation of a system at a multiple co-dimension bifurcation point. The authors then applied the method to the hydrodynamic problem of double diffusion. The approach has been strongly advocated by \citet{roberts14} and, along with \cite{roberts97}, has inspired the algorithmic procedure presented in \cite{carini15} for fluid flows. 

Subsequently, the work of \citet{cabre03a,cabre03b,cabre05,haro16} introduced the parametrization method which has been successfully utilized for the calculation of spectral sub-manifolds (SSM) \citep{jain22,vizzaccaro22,opreni23,thurnher24}. The framework is quite general and allows for any particular style of representation of the reduced dynamics. Indeed, the normal-form can be seen as a particular choice of parameterization. As a side note, the essential idea of imposing an arbitrary choice of parameterization to the invariant manifold and evaluating the corresponding asymptotic vectors was also expressed in \cite{roberts97}, albeit via more specific examples.

Often however, one is interested in systems slightly away from the bifurcation point, where the system parameters have been perturbed or small amplitude time dependent forcing has been applied. This leads us to the territory of non-autonomous dynamical systems. Theoretical work on invariant manifolds of general non-autonomous systems exists \citep{chicone97,siegmund02,hochs19,roberts22}, however, until fairly recently few practical applications to large systems had appeared. The applications to small systems often relied on integral transforms requiring convolutions in time \citep{cox91,siegmund02,roberts22}. Such operations tend to impose prohibitive computational costs for large systems. 

A simpler scenario exists when the time-dependence of the system can be specified with a finite (usually small) number of harmonics. The work presented in \cite{jain22,vizzaccaro22,thurnher24} leverages the parameterization method to develop a methodology for the computation of invariant manifolds of autonomous and non-autonomous mechanical systems which does not rely on convolutions in time. 
An alternate approach is to elevate parameter perturbations and harmonic forcings to state variables. This of course is not a novel idea and is routinely performed ``as a trick'' for small dynamical systems, with some examples of application to larger systems \citep{mercer90,cox91,jiang91,negi24,vizzaccaro24}. 
However, a thorough systematic examination of such an operation for general dynamical systems and its consequences for asymptotic evaluations of invariant manifolds is lacking. The current work investigates this in a general setting, illuminating the structure of the new tangent spaces and subsequent local extended center manifold. 

As it turns out this detailed examination reveals a few pleasant surprises. Despite a non-trivial expansion of the center subspace, high order asymptotic terms of the local manifold can be evaluated using the linear operators of the original (unextended) system. This is an important aspect for large scale systems since solution strategies for such systems often rely on exploiting the inherent structure of the system. Changing the system matrices due to the expansion of the state vector could potentially destroy the structure of the operators causing existing algorithms to fail. Fortunately, the system matrices require no modification. Secondly, since the external forcing and parameter perturbations are treated at par with internal modes, the algorithmic treatment is unified. A single algorithm therefore incorporates all the three cases of center manifold approximations - unperturbed system, parameter perturbations and harmonically forced systems. 

A more subtle but significant difference arises with the interpretation of the invariant manifold in the non-autonomous setting. The approach in \cite{jain22} is to consider the tori resulting from the deformation of the hyperbolic fixed point of the original system under small amplitude harmonic forcing, and evaluate the phase dependent invariant manifold attached to this tori. However, as pointed out in \cite{jain22}, the tori is not guaranteed to exist for a non-hyperbolic fixed point. A center manifold approximation is the appropriate setting in such a case. The tori (should it exist) is simply part of the larger center manifold emanating from the fixed point of the extended system, which is guaranteed to exist locally by the center manifold theorem. The investigation in the current work bears similarities to the work of \citet{vizzaccaro24} who also treat forced mechanical systems using an extension of the system however, they do not consider parameter extensions (and constant forcing terms). Additionally, the authors assume the original linearized system to be non-singular with damped eigenvalues, possibly to stay within the hyperbolic fixed point assumption although, this was not explicitly stated. However, parameter perturbations (and steady forcing) lead to complex interaction terms precisely in the scenario of a singular system and hence, require careful treatment. 

The current work therefore can be considered complementary to that of \cite{jain22} and \cite{vizzaccaro24}. It presents in a general setting the treatment of parameter perturbations, something that has been treated by way of an example of a perturbed Lorenz system in \cite{jain22} and avoided in \cite{vizzaccaro24}. In addition, it allows for the construction of center manifolds of non-autonomous systems even in the case of non-hyperbolic fixed points (where exact resonances with external frequencies and parameter perturbation may occur),
and unifies the algorithmic treatment for all cases.  

The remainder of the manuscript is laid out as follows. Section~\ref{sec:problem_setup} presents the problem in a general setting and subsequently evaluates the structure of the extended tangent spaces under parameter perturbation and harmonic forcings. For clarity, the tangent space extension for parameter perturbations and harmonic forcings will be treated as two distinct classes as they are usually understood in the literature. However, the algorithmic treatment makes it obvious that they can be subsumed into a single larger class through system extension. Section~\ref{sec:center_manifold_derivation} then builds the algorithms for the construction of asymptotic center manifolds of the extended systems while exploiting the tangent space structure derived earlier. The presented algorithm is then applied to a Ginzburg-Landau model under various conditions of system extension in Section~\ref{sec:ginzburg_landau}. Concluding remarks are made in Section~\ref{sec:discussion}.










